\section*{Hoofdstuk \ref{ch:b1:geological-maps} --- Geologische kaarten}

\begin{solution}{exo:b1:geological-maps:1}
\omterm{def:b1:geological-maps:attitude}{Strijk}: \omterm{def:b1:spherical-astronomy:horizontal}{azimut} van de horizontale lijnen van het vlak; \omterm{def:b1:geological-maps:attitude}{helling}: hoek tussen het vlak en de
horizontale, gemeten loodrecht op de \omterm{def:b1:geological-maps:attitude}{strijk}; \omterm{def:b1:geological-maps:attitude}{hellingrichting}: \omterm{def:b1:spherical-astronomy:horizontal}{azimut} van de steilste neerwaartse
lijn. Voor $135/40$ is de \omterm{def:b1:geological-maps:attitude}{strijk} volgens de rechterhandregel $135 - 90 = \ang{045}$. Symbool: een
strijklijn gericht N\,45\,O, met een streep aan de zuidoostkant (naar \ang{135}) met daarbij 40.
\end{solution}

\begin{solution}{exo:b1:geological-maps:2}
$z = 500 - \tan\ang{30}\,(x\sin\ang{90} + y\cos\ang{90}) = 500 - 0.577\,x$.
Op $(1000, 0)$: $500 - 577 = \qty{-77}{m}$. Op $(0, 1000)$: \qty{500}{m}, omdat het vlak
horizontaal is langs zijn noord--zuid \omterm{def:b1:geological-maps:attitude}{strijk}.
\end{solution}

\begin{solution}{exo:b1:geological-maps:3}
Een spoor dat een hoogtelijn volgt, is een horizontaal contact op \qty{350}{m}. Een spoor dat recht
over heuvels en dalen loopt, hoort bij een verticaal contact.
\end{solution}

\begin{solution}{exo:b1:geological-maps:4}
Langs de rivier toont de kern van de \omterm{def:b1:geological-maps:fold}{antiklinaal} de schist en de graniet (onder de zandsteen); de
kern van de \omterm{def:b1:geological-maps:fold}{synklinaal} wordt ingenomen door het krijt. De oudste eenheden liggen het laagst in de
reeks: waar de lagen omhoog buigen en de erosie er het diepst in snijdt, komen de oudste tevoorschijn;
waar zij doorzakken, blijven de jongste in het midden bewaard.
\end{solution}

\begin{solution}{exo:b1:geological-maps:5}
De eerste twee punten liggen op dezelfde hoogte: de lijn $y = 0$ is een strijklijn, dus de \omterm{def:b1:geological-maps:attitude}{strijk} is
oost--west. Het contact ligt \qty{100}{m} lager op \qty{500}{m} naar het noorden: het helt naar het noorden,
$\tan\delta = 100/500 = 0.2$, $\delta \approx \ang{11.3}$; \omterm{def:b1:geological-maps:attitude}{hellingrichting} \ang{000}, \omterm{def:b1:geological-maps:attitude}{strijk} \ang{270} volgens de
rechterhandregel. Met de stelling: $\overrightarrow{P_1P_2} = (1000, 0, 0)$,
$\overrightarrow{P_1P_3} = (0, 500, -100)$, $\vec n = (0, 100\,000,
500\,000)$, $\cos\delta = 500\,000/509\,902 \approx 0.981$, $\delta \approx
\ang{11.3}$, en $(n_x, n_y) = (0, 100\,000)$ wijst naar het noorden: hetzelfde resultaat.
\end{solution}

\begin{solution}{exo:b1:geological-maps:6}
Stroomafwaarts is het noorden, $s = 0.03$ (\ang{1.72}). (a) $a = 0 < s$: de V wijst stroomopwaarts
(zuid), langs de hoogtelijnen. (b) $a = -\tan\ang{20} < 0 < s$: de V wijst stroomopwaarts.
(c) $a = \tan\ang{1} \approx 0.017 < s$: de V wijst stroomopwaarts. (d) $a = \tan\ang{20}
\approx 0.364 > s$: de V wijst stroomafwaarts (noord).
\end{solution}

\begin{solution}{exo:b1:geological-maps:7}
$\xi_{\text{bodem}} - \xi_{\text{wand}} = h/(a - s) = 80/(0.364 - 0.03)
\approx \qty{240}{m}$: het spoor op de wand ligt ongeveer \qty{240}{m} stroomopwaarts van
de top.
\end{solution}

\begin{solution}{exo:b1:geological-maps:8}
De sluiting van de \omterm{def:b1:geological-maps:fold}{antiklinaal} wijst in de richting van de \omterm{def:b1:geological-maps:attitude}{helling} van de as: de \omterm{def:b1:geological-maps:fold}{plooias} helt naar het
noorden. Met een \omterm{def:b1:geological-maps:fold}{plooias} die naar het noorden helt, sluiten de \omterm{def:b1:geological-maps:geological-map}{ontsluitingen} van een \omterm{def:b1:geological-maps:fold}{synklinaal} zich
naar het zuiden.
\end{solution}

\begin{solution}{exo:b1:geological-maps:9}
De kant ten oosten-zuidoosten toont jongere eenheden (krijt en schalie tegenover oudere eenheden
er tegenover): dat blok is omlaag bewogen. De tikjes van het breuksymbool liggen aan die kant, de
omlaag bewogen kant.
\end{solution}

\begin{solution}{exo:b1:geological-maps:10}
Een lijn die beide vlakken gemeen hebben, staat loodrecht op beide normalen, dus is evenwijdig aan
$\vec n_1 \times \vec n_2$. Voor $090/30$: $\vec n_1 = (0.577, 0,
1)$; voor $270/30$: $\vec n_2 = (-0.577, 0, 1)$. Dan is $\vec n_1 \times \vec n_2
= (0\cdot 1 - 1\cdot 0,\ 1\cdot(-0.577) - 0.577\cdot 1,\ 0) = (0, -1.155,
0)$: een horizontale noord--zuidlijn. Het eerste vlak helt naar het oosten, het tweede naar het
westen: zij helt van de lijn af, zoals de kruin van een \omterm{def:b1:geological-maps:fold}{antiklinaal} met een horizontale as.
\end{solution}

\begin{solution}{exo:b1:geological-maps:11}
Het spoor van de verticale dijk is de rechte lijn door de oorsprong in de richting \ang{045}, op elke
hoogte: op de heuvel ligt het op diezelfde lijn. De hellende dijk is het vlak $z = -\tan\ang{60}\,d$, waarin
$d$ de horizontale afstand is tot de strijklijn door de oorsprong, positief naar het zuidoosten gerekend.
Op $z = \qty{300}{m}$: $d = -300/\tan\ang{60}
\approx \qty{-173}{m}$: het spoor ligt \qty{173}{m} noordwest van de strijklijn door de oorsprong,
aan de opwaartse kant van de \omterm{def:b1:geological-maps:attitude}{helling}.
\end{solution}

\begin{solution}{exo:b1:geological-maps:12}
\[
  \frac{\partial z}{\partial x} \approx \frac{405 - 420}{100} = -0.15, \qquad
  \frac{\partial z}{\partial y} \approx \frac{412 - 420}{100} = -0.08,
\]
dus de \omterm{def:b1:geological-maps:attitude}{helling} is $\arctan\sqrt{0.15^2 + 0.08^2} = \arctan 0.17 \approx
\ang{9.6}$. De grond daalt langs $-\nabla z = (0.15, 0.08)$, waarvan het \omterm{def:b1:spherical-astronomy:horizontal}{azimut}
$\arctan(0.15/0.08) \approx \ang{62}$ is: de cel is naar het oost-noordoosten gericht.
\end{solution}

\begin{solution}{pb:b1:geological-maps:1}
\textbf{1.} Schist, graniet, zandsteen, kalksteen, schalie, krijt, alluvium.

\textbf{2.} Van west naar oost: krijt en schalie, kalksteen, zandsteen, dan schist en graniet nabij
$x = \qty{3.5}{km}$, zandsteen, kalksteen en schalie opnieuw, dan krijt; alluvium over de hele bodem.

\textbf{3.} De zandsteen, met de schist en de graniet eronder langs de rivier: de oudste eenheden in
de kern, een \omterm{def:b1:geological-maps:fold}{antiklinaal}.

\textbf{4.} Het krijt, de jongste eenheid: een \omterm{def:b1:geological-maps:fold}{synklinaal}.

\textbf{5.} De basis van de zandsteen is het hoogst op de as van de \omterm{def:b1:geological-maps:fold}{antiklinaal}, en alleen langs de rivier
heeft de erosie er onder gesneden.

\textbf{6.} Een discordantie: de zandsteen rust op twee verschillende oudere eenheden, schist en graniet, en
snijdt de omtrek van de graniet af.

\textbf{7.} $435 - 261 = 174$, $432 - 428 = 4$, $512 - 586 = \qty{-74}{m}$.

\textbf{8.} $\overrightarrow{E_1E_2} = (600, 1200, -170)$,
$\overrightarrow{E_1E_3} = (1100, 400, -248)$.

\textbf{9.} Het uitwendig product is
\[
  (-229\,600,\ -38\,200,\ -1\,080\,000);
\]
naar boven georiënteerd, $\vec n_E = (229\,600,\ 38\,200,\ 1\,080\,000)$.

\textbf{10.} $\|\vec n_E\| \approx 1\,104\,800$; $\cos\delta \approx 0.9776$,
$\delta \approx \ang{12.2}$.

\textbf{11.} \omterm{def:b1:spherical-astronomy:horizontal}{Azimut} van $(229\,600, 38\,200)$: $\arctan(229\,600/38\,200)
\approx \ang{80.6}$; oriëntatie $081/12$, \omterm{def:b1:geological-maps:attitude}{strijk} $\approx \ang{351}$.

\textbf{12.} Ja: het symbool op de oostflank heeft een strijklijn nabij
noord--zuid en een helling-streep naar het oosten, met 12 erbij.

\textbf{13.} $524 - 669 = -145$, $421 - 459 = -38$, $319 - 209 =
\qty{110}{m}$.

\textbf{14.} $\overrightarrow{W_1W_2} = (700, 1200, 107)$ en
$\overrightarrow{W_1W_3} = (1100, -600, 255)$. Hun uitwendig product, naar boven
georiënteerd, is
\[
  (370\,200,\ -60\,800,\ -1\,740\,000) \quad\longrightarrow\quad
  \vec n_W = (-370\,200,\ 60\,800,\ 1\,740\,000).
\]

\textbf{15.} $\cos\delta = 1\,740\,000/\|\vec n_W\| \approx 0.9776$, $\delta
\approx \ang{12.2}$; \omterm{def:b1:geological-maps:attitude}{hellingrichting}: \omterm{def:b1:spherical-astronomy:horizontal}{azimut} van $(-370\,200, 60\,800)$,
$\approx \ang{279.3}$. Oriëntatie $279/12$, \omterm{def:b1:geological-maps:attitude}{strijk} $\approx \ang{189}$: de westflank helt naar het
west-noordwesten, de oostflank naar het oost-noordoosten.

\textbf{16.} $\vec n_E \times \vec n_W = (804 \times 10^6,\ -799\,320 \times
10^6,\ 28\,101 \times 10^6)$ (afgerond).

\textbf{17.} De neerwaartse richting is $(-804, 799\,320, -28\,101) \times
10^6$: richting $\arctan(-804/799\,320) \approx \ang{0}$, dat is het noorden; \omterm{def:b1:geological-maps:attitude}{helling} $\arctan\bigl(28\,101/\sqrt{804^2 + 799\,320^2}\bigr) \approx \ang{2.0}$. De
\omterm{def:b1:geological-maps:fold}{plooias} helt \ang{2} naar het noorden.

\textbf{18.} $8000 \times \tan\ang{2.0} \approx \qty{280}{m}$ lager, dus nabij
\qty{-50}{m}. De kern van de \omterm{def:b1:geological-maps:fold}{antiklinaal} is het hoogst nabij de rivier en zakt naar het noorden,
zodat haar \omterm{def:b1:geological-maps:geological-map}{ontsluiting} zich naar het noorden sluit, zoals de kaart toont.

\textbf{19.} $\overrightarrow{F_1F_2} = (1627, 2500, -276)$ en
$\overrightarrow{F_1F_3} = (3124, 5500, 77)$ geven de opwaartse normaal
\[
  (1\,710\,500,\ -987\,503,\ 1\,138\,500), \qquad
  \delta = \arccos\frac{1\,138\,500}{2\,279\,700} \approx \ang{60},
\]
en een \omterm{def:b1:geological-maps:attitude}{hellingrichting} $\approx \ang{120}$. Oriëntatie $120/60$, \omterm{def:b1:geological-maps:attitude}{strijk}
\ang{030}.

\textbf{20.} De breuk helt naar het oost-zuidoosten, dus die kant is het hangende blok, het blok dat boven
het hellende breukvlak ligt; het is omlaag bewogen (jongere eenheden aan het oppervlak): een normale breuk.

\textbf{21.} Het breukvlak helt oost-zuidoostwaarts, stroomafwaarts en veel steiler dan de bodem van de
rivier daalt: lager terrein komt verder in de richting van de \omterm{def:b1:geological-maps:attitude}{helling} op het vlak uit, zodat het spoor in de
vallei naar het oost-zuidoosten buigt --- de V-regel toegepast op het breukvlak.

\textbf{22.} Intrusie van het graniet; erosie vóór de zandsteen; afzetting van het krijt; plooiing;
breukvorming; afzetting van het alluvium.

\textbf{23.} Na 90 miljoen jaar geleden (het krijt is geplooid) en vóór 2.58 miljoen jaar geleden (het
alluvium is niet gebroken); de plooiing kwam eerst, de breukvorming daarna.

\textbf{24.} De breuk daalt \qty{1405}{m} van F1 (\qty{405}{m}) tot
\qty{-1000}{m}; bij een \omterm{def:b1:geological-maps:attitude}{helling} van \ang{60} is dat een horizontale afstand van
$1405/\tan\ang{60} \approx \qty{811}{m}$ in de \omterm{def:b1:geological-maps:attitude}{helling}. De put moet ongeveer
\textbf{\qty[text-series-to-math=true]{810}{m} oost-zuidoost van F1} worden geboord (\omterm{def:b1:spherical-astronomy:horizontal}{azimut} \ang{120}).
\end{solution}
